\subsection{Detection and classification (Part B)}
\label{sec:detection-methods}
We find peaks with a prominence above the threshold $k\sigma$ (Section~\ref{sec:threshold})
on both channels. One tissue piece often gives
several peaks, so peaks less than \SI{1}{\second} apart are grouped into one \emph{event}. An
event represents either a piece of tissue or water. Each event is labelled from its largest peak.

Brightness comes first: an event whose largest peak is above the no-dye level is
\texttt{dyed}. The no-dye level (68.9 counts) is the top 0.1\,\% of the peaks in the water-only
recordings.

A dim event is either tissue without dye or water, which we separate with a two-channel check.
Solid tissue also gives a peak on \texttt{adc1}, water usually does not. The check passes if an
\texttt{adc1} peak of at least 5\,\% of the height lies within \SI{\pm 20}{\milli\second} of the
largest \texttt{adc2} peak, after correcting the delay between the channels (measured for each
recording). The event is then \texttt{solid\_no\_dye}, otherwise \texttt{water}. If \texttt{adc1}
hardly responds (Part A), or if the check passes by chance in more than 20\,\% of the cases (see
below), it is not used and dim events are \texttt{undecided}.

Besides the known answers, we used four measures to check the detection. The \emph{chance rate}
is the share of events that pass the two-channel check when \texttt{adc1} is shifted by
2--\SI{20}{\second}, so that real pairs cannot match. The \emph{confirmation rate} is the share
of \texttt{dyed} pieces that also pass the check; dyed pieces are solid, so it should be high.
The \emph{noisiness} is the total up-and-down movement of \texttt{adc2} inside an event divided
by twice its largest peak (1 for one clean hump). The \emph{signal outside the events} is the
share of the signal above $3\sigma$ in a suction that is not inside any event; it grows if small
peaks are lost under the threshold.
